% q0016.tex — To be or Not to be (Independent) — Phillips Curve / CB Independence
\question{
  id        = {0016},
  title     = {To Be or Not to Be (Independent)},
  source    = {OBECON},
  year      = {2026},
  round     = {Seletiva IEO -- Phase A},
  language  = {English},
  topic     = {Macro},
  subtopic  = {Monetary Policy / Phillips Curve},
  type      = {objective},
  statement = {%
    The image below shows reports that President Trump once considered firing
    Federal Reserve Chair Jerome Powell over disagreements about monetary
    policy. Such an event highlights the tension between political goals and
    central bank independence, especially when it comes to controlling
    inflation and unemployment.

    \begin{tcolorbox}[
      colback=white, colframe=black!40, arc=2pt, boxrule=0.5pt,
      left=6pt, right=6pt, top=4pt, bottom=4pt,
    ]
      {\large\bfseries\itshape Trump Has Draft of Letter to Fire Fed Chair.
      He Asked Republicans if He Should Send It.}\\[4pt]
      {\small\textit{The president waved a copy of a draft letter firing
      Jerome H.\ Powell at a meeting in the Oval Office with House
      Republicans. It remains to be seen whether he follows through
      with his threat.}}
    \end{tcolorbox}

    \vspace{6pt}
    Imagine, then, that you're analyzing the economic implications of such a
    standoff. Suppose the economy has a natural unemployment rate of 5\% and
    that the short-run Phillips curve is given by:
    \[
      \pi = \pi^e - 0.5\,(u - u_N)
    \]
    where $\pi$ is the actual inflation, $\pi^e$ is the expected inflation,
    $u$ is the unemployment rate and $u_N$ is the natural rate of
    unemployment. At first the economy has an unemployment rate of 3\% and
    expected inflation of 2\%. After the announcement, inflation expectations
    rose to 4\%. Suppose the Fed maintains its current policy stance. What
    are the implications for inflation and unemployment in the short run and
    long run?
  },
  choices   = {%
    \item In the short run, inflation increases to 5\% while unemployment
          remains at 3\%; in the long run, unemployment returns to the
          natural rate of 5\% and inflation stabilizes at 4\%.
    \item In the short run, inflation decreases to 5\% while unemployment
          remains at 3\%; in the long run, unemployment returns to the
          natural rate of 5\% and inflation stabilizes at 4\%.
    \item In the short run, inflation increases to 5\% while unemployment
          remains at 3\%; in the long run, unemployment returns to the
          natural rate of 5\% and inflation stabilizes at 3\%.
    \item In the short run, inflation decreases to 5\% while unemployment
          remains at 3\%; in the long run, unemployment returns to the
          natural rate of 3\% and inflation stabilizes at 3\%.
    \item In the short run, inflation remains at 3\% while unemployment
          increases to 5\%; in the long run, unemployment returns to the
          natural rate of 3\% and inflation stabilizes at 4\%.
  },
  answer    = {},
  solution  = {%
    % TODO: add solution
  },
}
